% Part: first-order-logic
% Chapter: axiomatic-deduction
% Section: proof-theoretic-notions

\documentclass[../../../include/open-logic-section]{subfiles}

\begin{document}

\iftag{FOL}
      {\olfileid{fol}{axd}{ptn}}
      {\olfileid{pl}{axd}{ptn}}

\olsection{સાબિતી-સૈદ્ધાંતિક ખ્યાલો}

\begin{explain}
આપણે અનેક મહત્વના અર્થાનુસારી ખ્યાલો
(\iftag{FOL}{પ્રામાણ્ય}{પુનરુક્તિ}, ફલિતતા અને સંતોષ્યતા) વ્યાખ્યાયિત
કર્યા છે; એ જ રીતે હવે તેમને અનુરૂપ \emph{સાબિતી-સૈદ્ધાંતિક ખ્યાલો}
વ્યાખ્યાયિત કરીએ છીએ. આ ખ્યાલો !!{structure}sમાં !!{sentence}sના
સંતોષનો આશરો લઈને નહિ, પરંતુ અમુક સૂત્રોની !!{derivability} અથવા
!!{nonderivability}નો આશરો લઈને વ્યાખ્યાયિત થાય છે. આ બંને પ્રકારના
ખ્યાલો મળી જાય છે એ એક મહત્વની શોધ હતી. તેઓ મળે છે તે જ
\emph{યથાર્થતા} અને \emph{પૂર્ણતાના પ્રમેયો}નો વિષય છે.
\end{explain}

\begin{defn}[!!^{derivability}]
જો~$\Gamma$માંથી મળતી એવી !!a{derivation} હોય જે~$!A$માં પૂરી થાય,
તો !!^a{formula}~$!A$ એ~$\Gamma$માંથી \emph{!!{derivable}} છે; તેને
$\Gamma \Proves !A$ લખીએ છીએ.
\end{defn}

\begin{defn}[પ્રમેયો]
ખાલી ગણમાંથી~$!A$ની !!a{derivation} હોય, તો !!^a{formula}~$!A$
\emph{પ્રમેય} છે. $!A$ પ્રમેય હોય તો $\Proves !A$ અને ન હોય તો
$\Proves/ !A$ લખીએ છીએ.
\end{defn}

\begin{defn}[સુસંગતતા]
!!{formula}sનો ગણ~$\Gamma$ \emph{સુસંગત} છે જો અને માત્ર જો
$\Gamma\Proves/ \lfalse$; નહિ તો તે \emph{અસુસંગત} છે.
\end{defn}

\begin{prop}[સ્વવાચકતા]
\ollabel{prop:reflexivity}
જો $!A \in \Gamma$ હોય, તો $\Gamma \Proves !A$.
\end{prop}

\begin{proof}
  એકલું !!{formula}~$!A$ એ~$\Gamma$માંથી $!A$ની !!a{derivation} છે.
\end{proof}

\begin{prop}[એકદિશવર્ધિતા]
\ollabel{prop:monotonicity}
જો $\Gamma \subseteq \Delta$ અને $\Gamma \Proves !A$ હોય, તો
$\Delta \Proves !A$.
\end{prop}

\begin{proof}
$\Gamma$માંથી $!A$ની કોઈ પણ !!{derivation} એ~$\Delta$માંથી $!A$ની પણ
!!a{derivation} છે.
\end{proof}

\begin{prop}[પરંપરિતતા]
\ollabel{prop:transitivity}
જો $\Gamma \Proves !A$ અને $\{!A\} \cup \Delta \Proves !B$ હોય, તો
$\Gamma \cup \Delta \Proves !B$.
\end{prop}

\begin{proof}
  ધારો કે $\{!A\} \cup \Delta \Proves !B$. તેથી~$\{!A\} \cup
  \Delta$માંથી એવી !!a{derivation} $!B_1$, \dots, $!B_l = !B$ છે.
  એ !!{derivation}નાં અમુક પગલાં યોગ્ય હશે કારણ કે તેમનું પ્રમાણન
  અગાઉની કોઈ પંક્તિ~$!B_i = !A$નો સંદર્ભ આપતા નિયમથી થાય છે. પૂર્વધારણા
  મુજબ~$\Gamma$માંથી $!A$ની !!a{derivation} છે, એટલે એવી
  !!a{derivation}~$!A_1$, \dots, $!A_k = !A$ છે જેમાં દરેક $!A_i$ કાં
  તો સ્વયંસિદ્ધિ, કાં~$\Gamma$નું !!a{element}, અથવા અનુમાનના નિયમ વડે
  યોગ્ય છે. હવે નીચેની શ્રેણી વિચારો:
  \[
  !A_1, \dots, !A_k = !A, !B_1, \dots, !B_l = !B.
  \]
  આ~$\Gamma \cup \Delta$માંથી $!B$ની યોગ્ય !!{derivation} છે, કારણ કે
  દરેક $!B_i = !A$ હવે એ જ કારણથી પ્રમાણિત થાય છે જે~$!A_k = !A$ને
  પ્રમાણિત કરે છે.
\end{proof}

\sourcecorrection{OLAX-004}{મૂળની
\texttt{proof-theoretic-notions.tex}, પંક્તિઓ 82--83માં દરેક
$!B_i=!A$ને એ જ ``નિયમ'' પ્રમાણિત કરે છે જે $!A_k=!A$ને પ્રમાણિત કરે છે
એમ લખ્યું છે. $!A_k$ સ્વયંસિદ્ધિ અથવા~$\Gamma$નું ઘટક પણ હોઈ શકે છે,
તેથી અહીં વધુ ચોક્કસ ``એ જ કારણ'' લખ્યું છે.}

\sourcecorrection{OLAX-016}{મૂળની
\texttt{proof-theoretic-notions.tex}, પંક્તિ~82માં સૂત્ર-ચલ~$B_i$
પહેલાંનો આંતરિક \texttt{!} સંકેત ખૂટે છે, જ્યારે આ જ નિષ્પત્તિની
અગાઉની બધી પંક્તિઓમાં $!B_i$ લખાયું છે. અહીં એ સંકેત પુનઃસ્થાપિત
કર્યો છે.}

નોંધો કે ખાસ કરીને તેનો અર્થ એવો થાય છે કે જો $\Gamma \Proves !A$ અને
$!A \Proves !B$ હોય, તો $\Gamma \Proves !B$. વળી, જો
$!A_1, \dots, !A_n \Proves !B$ અને દરેક~$i$ માટે
$\Gamma \Proves !A_i$ હોય, તો $\Gamma \Proves !B$ એવું પણ ફલિત થાય છે.

\begin{prop}
\ollabel{prop:incons}
$\Gamma$ અસુસંગત છે જો અને માત્ર જો દરેક~$!A$ માટે
$\Gamma \Proves !A$.
\end{prop}

\begin{proof}
પ્રશ્ન.
\end{proof}

\tagprob{FOL}
\begin{prob}
\olref[fol][axd][ptn]{prop:incons} સાબિત કરો.
\end{prob}
\tagendprob

\tagprob{notFOL}
\begin{prob}
\olref[pl][axd][ptn]{prop:incons} સાબિત કરો.
\end{prob}
\tagendprob

\begin{prop}[સઘનતા]
\ollabel{prop:proves-compact}
  \begin{enumerate}
  \item જો $\Gamma \Proves !A$, તો એવો સાન્ત ઉપગણ $\Gamma_0
    \subseteq \Gamma$ છે કે $\Gamma_0 \Proves !A$.
  \item જો~$\Gamma$નો દરેક સાન્ત ઉપગણ સુસંગત હોય, તો $\Gamma$
    સુસંગત છે.
  \end{enumerate}
\end{prop}

\begin{proof}
  \begin{enumerate}
    \item જો $\Gamma \Proves !A$, તો !!{formula}sની એવી સાન્ત શ્રેણી
      $!A_1$, \dots,~$!A_n$ છે કે $!A \ident !A_n$ અને દરેક $!A_i$
      કાં તો તાર્કિક સ્વયંસિદ્ધિ, કાં~$\Gamma$નું !!a{element}, અથવા
      અગાઉનાં !!{formula}sમાંથી મોડસ પોનેન્સ વડે મળતું સૂત્ર છે.
      $\Gamma_0$ને એવા $!A_i$નો ગણ લો જે~$\Gamma$માં છે. ત્યારે એ
      !!{derivation} એ~$\Gamma_0$માંથી પણ !!a{derivation} છે અને તેથી
      $\Gamma_0 \Proves !A$.
    \item $!A \ident \lfalse$ના વિશેષ કિસ્સામાં આ~(1)નું
      પ્રતિપક્ષી વિધાન છે.
\end{enumerate}
\end{proof}

\end{document}
